\honest{Guessing the Teacher's Password}{Eliezer Yudkowsky}{2007}

When I was nine I read Feynman's \textsc{QED} and knew that light, sound and matter were waves. \nb{\textsc{QED} tells its readers that light ``behaves like particles,'' and warns that at school they were ``probably told something about light behaving like waves.'' The post does not say where the young author's ``waves'' came from.} Later, reading the \textsc{Feynman Lectures on Physics}, I thought about the wave equation for three days until it was obvious, and realized I had never known what ``wave'' meant to a physicist.

We tend to think that a student who answers ``Waves!'' has made a true statement, since we accept the same word from a physicist. That is a bad habit learned at school. Words have no intrinsic definitions. If I hear ``beaver'' and think of a large rodent, that is a fact about my mind, not about the syllables. ``Made of waves'' can call up a hypothesis in someone's mind, but the syllables themselves are neither right nor wrong. School gives gold stars for saying them, so students come to think that what they say has a truth-value of its own. \nb{On a standard view in philosophy of language, from Hilary Putnam and Tyler Burge, a speaker who uses a word as the experts do, and would accept their correction, means what the experts mean. On that view the student's answer says what the physicist's says, and what the student lacks is understanding.}

Ask a student why the far side of a metal plate is warmer. ``I don't know'' earns no gold star, so the student offers one of the teacher's phrases: ``Eh, maybe because of heat conduction?'' This is not a hypothesis about the plate. It is ``an attempt to guess the teacher's password.'' Unlearning the habit starts with noticing the difference between an explanation and a password. \nb{Feynman described the same thing in \textsc{Surely You're Joking, Mr. Feynman!}: Brazilian students who could define Brewster's angle but did not recognize it in the light off the bay had ``memorized everything, but they didn't know what anything meant.''}

An equation connects to experience only if you use it to predict the next measurement. Even a student who pictures heat flowing and holds a match to the plate is connected. Yet if you are not using the diffusion equation, the connection ``is severed as though by a knife.'' \nb{Read literally, that includes the student with the match.}

Is this too harsh? Can't ``heat conduction?'' be a first step? Maybe, if you do not think you are hunting for a password. But with no teacher to tell you that you failed, you may think ``Light is wakalixes'' a good explanation. \nb{The word is Feynman's.} It happened to me at nine, ``not because I was stupid, but because this is what happens by default.'' ``Humanity stayed stuck in holes like this for thousands of years.'' I give no historical example.

Maybe we could drill students that only anticipation-controllers count. Then ``heat conduction?'' could lead to questions: what the phrase means, what the equation predicts, whether the near side only feels cooler, how to watch heat spread with a match.

Once more, in italics: this ``happened to the whole human species for thousands of years.''
